% !TeX root = ../main.tex
% 第 4 章 非整数维数集合的结构 (对应原书 pp. 54--63)

\chapter{非整数维数集合的结构}\label{chap:4}

\section{引言}

本章考察 $s$ 非整数时 $\mathbb{R}^n$ 中 $s$-集的局部性质.基本结果是：任何这样的集合都是非正则的，即它的下圆形密度或下球面密度在几乎所有的点处严格小于 $1$.事实上，更强的结果——它的密度在几乎所有的点处不存在——也已被建立.与之前一样，我们还考察适当定义的切线的存在性，并证明存在这种切线的点之集必有零测度.

对平面的子集而言，这一工作完全归功于 Marstrand（\citeyear{marstrand1954a,marstrand1955}），前一篇文章给出了非常完整的论述.与整数维数的集合一样，高维类似物呈现出巨大的困难；其自然的推广最终由 \textcite{marstrand1964} 证明.

\section{\texorpdfstring{$0<s<1$ 的 $s$-集}{0<s<1 的 s-集}}

首先考虑 $\mathbb{R}^n$ 中 $s$ 严格小于 $1$ 的 \newterm{s-set-lt1}{$s$-集}{$s$-set (s < 1)}.此时基本性质（包括密度几乎处处不存在）相对容易得到.

关于这类集合的下述拓扑观察有时是有用的.

\begin{lemma}\label{lem:4.1}
$\mathbb{R}^n$ 中满足 $0<s<1$ 的 $s$-集 $E$ 是完全不连通的.
\end{lemma}

\begin{proof}
设 $x$ 与 $y$ 是 $E$ 的同一连通分支中两个不同的点.定义映射 $f:\mathbb{R}^n\to[0,\infty)$ 为 $f(z)=|z-x|$.由于 $f$ 不增大距离，由\cref{lem:1.8} 得 $\mathcal{H}^s(f(E))\leq\mathcal{H}^s(E)<\infty$.由于 $s<1$，$f(E)$ 是 $\mathbb{R}$ 的勒贝格测度为零的子集，特别地其补集稠密.取数 $r$ 使 $r\notin f(E)$ 且 $0<r<f(y)$，则
\[
  E=\{z\in E:|z-x|<r\}\cup\{z\in E:|z-x|>r\}.
\]
这是 $E$ 的一个开分解，$x$ 属于其中一部分，$y$ 属于另一部分，与 $x,y$ 属于 $E$ 的同一连通分支的假设矛盾.
\end{proof}

下述密度性质的证明是 \textcite{marstrand1954a} 的证明的一个略微缩短的版本.

\begin{theorem}\label{thm:4.2}
若 $E$ 是满足 $0<s<1$ 的 $s$-集，则密度 $D^s(E,x)$ 在 $E$ 的几乎每一点处不存在.
\end{theorem}

\begin{proof}
若结论不成立，则由\cref{cor:2.5}，$E$ 有一个正测度的可测子集，在其上密度存在且至少为 $2^{-s}>\frac{1}{2}$.选取足够小的 $\rho$，可找到 $s$-集 $F\subset E$ 使得当 $x\in F$ 时 $D^s(E,x)$ 存在且
\begin{equation}\label{eq:4.1}
  \mathcal{H}^{s}\bigl(E\cap B_r(x)\bigr)>\frac{1}{2}(2r)^{s}
\end{equation}
对一切 $r\leq\rho$ 成立.由 $\mathcal{H}^s$ 的正则性还可设 $F$ 是闭集.设 $y$ 为 $F$ 的一个聚点，并取数 $\eta$ 使 $0<\eta<1$.若 $A_{r,\eta}$ 表示环形区域 $B_{r(1+\eta)}(y)\setminus B_{r(1-\eta)}(y)$，则当 $r\to 0$ 时
\begin{equation}\label{eq:4.2}
  \begin{aligned}
    (2r)^{-s}\mathcal{H}^{s}(E\cap A_{r,\eta})
    &=(2r)^{-s}\mathcal{H}^{s}\bigl(E\cap B_{r(1+\eta)}(y)\bigr)
    -(2r)^{-s}\mathcal{H}^{s}\bigl(E\cap B_{r(1-\eta)}(y)\bigr)\\
    &\to D^{s}(E,y)\bigl((1+\eta)^{s}-(1-\eta)^{s}\bigr).
  \end{aligned}
\end{equation}
另一方面，对任意小的 $r$ 值可找到 $x\in F$ 使 $|x-y|=r$.于是
\[
  B_{\frac{1}{2}r\eta}(x)\subset A_{r,\eta},
\]
故由 \eqref{eq:4.1}，对任意小的 $r$ 值（见图~\ref{fig:4.1}）有
\[
  \frac{1}{2}r^{s}\eta^{s}<\mathcal{H}^{s}\bigl(E\cap B_{\frac{1}{2}r\eta}(x)\bigr)\leq\mathcal{H}^{s}(E\cap A_{r,\eta}).
\]
利用 \eqref{eq:4.2} 得出
\[
  2^{-(s+1)}\eta^{s}\leq D^{s}(E,y)\bigl((1+\eta)^{s}-(1-\eta)^{s}\bigr)=D^{s}(E,y)\bigl(2s\eta+O(\eta^{2})\bigr)
\]
（$\eta\to 0$）.当 $s<1$ 时这是不可能的，定理由这一矛盾得证.
\end{proof}

\begin{figure}[htbp]
  \centering
  \includegraphics[width=0.34\textwidth]{fig4-1.png}
  \caption{\cref{thm:4.2} 证明中的环形区域}
  \label{fig:4.1}
\end{figure}

\begin{corollary}\label{cor:4.3}
任何满足 $0<s<1$ 的 $s$-集都是非正则的.
\end{corollary}

\textcite{marstrand1954a} 关于下角密度的下述结果，其证明在某些方面与\cref{thm:4.2} 的证明相仿.

\begin{theorem}\label{thm:4.4}
设 $\theta$ 为单位向量，$\phi<\frac{1}{2}\pi$.则若 $E$ 是满足 $0<s<1$ 的 $s$-集，在几乎所有的 $x\in E$ 处有 $D^s(E,x,\theta,\phi)=0$.
\end{theorem}

\begin{proof}
反设不然.则按通常方式可找到 $\rho_0$、$\alpha>0$ 与闭 $s$-集 $F\subset E$，使得当 $x\in F$ 且 $r\leq\rho_0$ 时
\begin{equation}\label{eq:4.3}
  \mathcal{H}^{s}\bigl(E\cap S_r(x,\theta,\phi)\bigr)>\alpha r^{s}.
\end{equation}
由\cref{cor:2.4,cor:2.5} 当然可以选取 $y\in F$ 使 $D^s(E\setminus F,y)=0$ 且 $\bar{D}^s(E,y)=c2^{-s}$，其中 $0<c<\infty$.于是，给定 $\varepsilon>0$，存在 $\rho_1\leq\rho_0$ 使当 $r\leq\rho_1$ 时
\begin{equation}\label{eq:4.4}
  \mathcal{H}^{s}\bigl((E\setminus F)\cap B_r(y)\bigr)<\varepsilon r^{s}
\end{equation}
且
\begin{equation}\label{eq:4.5}
  \mathcal{H}^{s}\bigl(E\cap B_r(y)\bigr)<(c+\varepsilon)r^{s}.
\end{equation}
另外，选取 $\rho\leq\frac{1}{2}\rho_1$ 使
\begin{equation}\label{eq:4.6}
  (c-\varepsilon)\rho^{s}<\mathcal{H}^{s}\bigl(E\cap B_\rho(y)\bigr).
\end{equation}

\begin{figure}[htbp]
  \centering
  \includegraphics[width=0.36\textwidth]{fig4-2.png}
  \caption{\cref{thm:4.4} 证明中的点 $x$ 的选取}
  \label{fig:4.2}
\end{figure}

设 $x$ 为 $F\cap B_\rho(y)$ 中使数量积 $x\cdot\theta$ 最大的点，见图~\ref{fig:4.2}.则
\[
  \bigl(F\cap B_\rho(y)\bigr)\cup\bigl(F\cap S_r(x,\theta,\phi)\bigr)\subset F\cap B_{\rho+r}(y),
\]
并且这个并除点 $x$ 外是不交的.于是
\[
  \mathcal{H}^{s}\bigl(F\cap B_\rho(y)\bigr)+\mathcal{H}^{s}\bigl(F\cap S_r(x,\theta,\phi)\bigr)\leq\mathcal{H}^{s}\bigl(F\cap B_{\rho+r}(y)\bigr),
\]
故
\[
  \mathcal{H}^{s}\bigl(E\cap B_\rho(y)\bigr)+\mathcal{H}^{s}\bigl(E\cap S_r(x,\theta,\phi)\bigr)
  \leq\mathcal{H}^{s}\bigl(E\cap B_{\rho+r}(y)\bigr)+2\mathcal{H}^{s}\bigl((E\setminus F)\cap B_{\rho+r}(y)\bigr).
\]
若 $r\leq\rho$，则 $\rho+r\leq\rho_1$，于是可用 \eqref{eq:4.3}--\eqref{eq:4.6} 得到
\[
  (c-\varepsilon)\rho^{s}+\alpha r^{s}<(c+\varepsilon)(\rho+r)^{s}+2\varepsilon(\rho+r)^{s}.
\]
令 $\gamma=r/\rho$，上式成为
\[
  c-\varepsilon+\alpha\gamma^{s}<(c+3\varepsilon)(1+\gamma)^{s},
\]
对 $0<\gamma<1$ 成立.这对一切 $\varepsilon>0$ 成立，故
\[
  c+\alpha\gamma^{s}\leq c(1+\gamma)^{s}\leq c+cs\gamma
\]
对 $0<\gamma<1$ 成立.这与 $\alpha$ 的正性矛盾.
\end{proof}

若取 $S_r\bigl(x,\theta,\frac{1}{2}\pi\bigr)$ 为开半圆来计算密度，则该结果对 $\phi=\frac{1}{2}\pi$ 仍然成立.Davies 曾观察到，这可作为 $\mathbb{R}$ 中单侧密度结果的直接 $n$ 维类似物得到；后者由 \textcite{besicovitch1968} 在他的倒数第二篇论文中给出了漂亮的证明.

$0<s<1$ 的 $s$-集的切线问题没有特别的意义，因为这些集合如此稀疏，以至于用线段作逼近的任何想法都变得相当无意义.

在一种意义下，切线在所有点处存在是可能的.设 $\Gamma$ 是 $\mathbb{R}^n$ 中由光滑双射 $\psi:\mathbb{R}\to\mathbb{R}^n$ 定义的光滑曲线，并设 $E$ 是 $\mathbb{R}$ 中的 $s$-集.容易看出 $\psi(E)$ 是 $\mathbb{R}^n$ 中的 $s$-集，而且对几乎所有的 $x\in\psi(E)$，$\Gamma$ 在 $x$ 处的切线就是 $\psi(E)$ 在 $x$ 处按 \eqref{eq:3.4} 意义的切线.

\section{\texorpdfstring{$s>1$ 的 $s$-集}{s>1 的 s-集}}

当 $s>1$ 时，\newterm{s-set-gt1}{$s$-集}{$s$-set (s > 1)}的密度与切线性质要复杂得多，我们把详细的论述限于平面的子集.

先处理切线问题.一种途径是重复 Besicovitch 对\cref{thm:3.29} 的证明，推出平面 $s$-集（$1<s<2$）在任意一对对顶角中上角密度之和几乎处处严格为正.Besicovitch 的证明依赖于所论集合与任何可求长曲线的交测度为零；由\cref{lem:3.2} 之后的注记，当 $s>1$ 时这当然成立.于是由切线的定义 \eqref{eq:3.4} 可得：满足 $1<s<2$ 的平面 $s$-集几乎处处没有切线.然而，更容易证明一个更强的结果，而且在涉及圆形密度的问题时，我们需要这种关于切向行为的更细致的信息.

遵循 \textcite{marstrand1954a}，我们定义弱切线.若 $s$-集 $E$ 满足 $D^s(x,E)>0$，且对每个 $\phi>0$，
\begin{equation}\label{eq:4.7}
  \varliminf_{r\to 0}r^{-s}\mathcal{H}^{s}\bigl(E\cap(B_r(x)\setminus S_r(x,\theta,\phi)\setminus S_r(x,-\theta,\phi))\bigr)=0,
\end{equation}
则称 $E$ 在 $x$ 处有方向 $\theta$ 的\newterm{weak tangent}{弱切线}{weak tangent}.（回忆 $S_r(x,\theta,\phi)$ 是由 $B_r(x)$ 中在 $x$ 处与从 $x$ 出发方向 $\theta$ 的半直线夹角至多为 $\phi$ 的那些点构成的扇形.）

与 \eqref{eq:3.4} 定义的切线不同，一个集合在一点处在原则上可以有许多弱切线.然而这种情形通常不会发生.

\begin{lemma}\label{lem:4.5}
设 $E$ 是 $\mathbb{R}^2$ 中满足 $1<s<2$ 的 $s$-集.则对几乎所有的 $x\in E$，
\[
  \bar{D}^{s}(E,x,\theta,\phi)\leq 4\cdot 10^{s}\phi^{s-1}
\]
对一切 $\theta$ 与一切 $\phi\leq\frac{1}{2}\pi$ 成立.
\end{lemma}

\begin{proof}
固定 $\rho>0$，令
\begin{equation}\label{eq:4.8}
  F=\{x\in E:\mathcal{H}^{s}\bigl(E\cap B_r(x)\bigr)<2^{s+1}r^{s}\ \text{对一切}\ r\leq\rho\}.
\end{equation}
取 $x\in F$ 及任意满足 $0<\phi\leq\frac{1}{2}\pi$ 的 $\theta$ 与 $\phi$.对 $i=1,2,\ldots$，令 $A_i$ 为圆环与扇形之交
\[
  A_i=S_{ir\phi}(x,\theta,\phi)\setminus S_{(i-1)r\phi}(x,\theta,\phi),
\]
则 $S_r(x,\theta,\phi)\subset\bigcup_{1}^{m}A_i\cup\{x\}$，其中整数 $m<2/\phi$.当 $r<\rho/20$ 时，每个 $A_i$ 的直径至多为 $10r\phi<\rho$，于是对每个含 $F$ 的点的 $A_i$ 应用 \eqref{eq:4.8} 并求和，
\[
  \mathcal{H}^{s}\bigl(F\cap S_r(x,\theta,\phi)\bigr)\leq 2\phi^{-1}2^{s+1}(10r\phi)^{s},
\]
即当 $r<\rho/20$ 时
\[
  (2r)^{-s}\mathcal{H}^{s}\bigl(F\cap S_r(x,\theta,\phi)\bigr)\leq 4\cdot 10^{s}\phi^{s-1}.
\]
于是在几乎所有的 $x\in F$ 处 $\bar{D}^s(F,x,\theta,\phi)\leq 4\cdot 10^{s}\phi^{s-1}$；又由\cref{cor:2.6}，在几乎所有的 $x\in F$ 处 $\bar{D}^s(E\setminus F,x)=0$，故引理结论在几乎所有的 $x\in F$ 处成立.由\cref{cor:2.5}，$E$ 的几乎所有的点对某个 $\rho>0$ 属于这样的 $F$，故结论得证.
\end{proof}

\begin{corollary}\label{cor:4.6}
若 $E$ 是 $\mathbb{R}^2$ 中满足 $1<s<2$ 的 $s$-集，则在 $E$ 的几乎所有的点处不存在弱切线.
\end{corollary}

\begin{proof}
由\cref{lem:4.5}，对几乎所有的 $x\in E$ 及一切 $\theta$ 与 $\phi<\frac{1}{2}\pi$，
\begin{align*}
  \varliminf_{r\to 0}(2r)^{-s}\mathcal{H}^{s}&\bigl(E\cap(B_r(x)\setminus S_r(x,\theta,\phi)\setminus S_r(x,-\theta,\phi))\bigr)\\
  &\geq\underline{D}^{s}(E,x)-\bar{D}^{s}(E,x,\theta,\phi)-\bar{D}^{s}(E,x,-\theta,\phi)
  \geq\underline{D}^{s}(E,x)-8\cdot 10^{s}\phi^{s-1}.
\end{align*}
在这样的点处，要么 $D^s(E,x)=0$，要么当 $\phi$ 充分小时 \eqref{eq:4.7} 不成立，故不存在弱切线.
\end{proof}

为证明 $1<s<2$ 时 $s$-集是非正则的，我们先考虑角密度而非圆形密度.

\begin{theorem}\label{thm:4.7}
设 $E$ 是 $\mathbb{R}^2$ 中满足 $1<s<2$ 的 $s$-集.则当 $\phi<\frac{1}{2}\pi$ 时，对几乎所有的 $x\in E$，下角密度 $D^s(E,x,\theta,\phi)$ 对某个 $\theta$ 为零.
\end{theorem}

\begin{proof}
固定正数 $\alpha$ 与 $\rho$，令
\begin{equation}\label{eq:4.9}
  F_0=\bigl\{x:\mathcal{H}^{s}\bigl(E\cap S_r(x,\theta,\phi)\bigr)>\alpha r^{s}
  \ \text{对一切}\ r\leq\rho\ \text{与一切}\ \theta\bigr\}.
\end{equation}
我们证明 $\mathcal{H}^s(F_0)=0$.假设不然，则由\cref{cor:2.5} 可找到 $\rho_1\leq\rho$ 与正测度集 $F\subset F_0$，使得当 $x\in F$ 且 $r\leq\rho_1$ 时
\begin{equation}\label{eq:4.10}
  \mathcal{H}^{s}\bigl(E\cap B_r(x)\bigr)<2^{s+1}r^{s}.
\end{equation}
由 $\mathcal{H}^s$ 的正则性还可设 $F$ 为闭集.

设 $y$ 为 $F$ 中使 $E\setminus F$ 的圆形密度为零的一点；由\cref{cor:2.6}，$F$ 的几乎任一点都可以.于是给定 $\varepsilon>0$，可找到 $\rho_2\leq\rho_1$ 使当 $r\leq\rho_2$ 时
\begin{equation}\label{eq:4.11}
  \mathcal{H}^{s}\bigl((E\setminus F)\cap B_r(y)\bigr)<\varepsilon r^{s}.
\end{equation}
现在我们在圆盘 $B_{\rho_2}(y)$ 内工作.首先证明 $B_{\frac{1}{2}\rho_2}(y)$ 中有相对远离 $F$ 的点.设对某个 $\gamma\leq\frac{1}{2}\rho_2$，$B_{\frac{1}{2}\rho_2}(y)$ 的所有点与 $F$ 的距离都不超过 $\gamma$.则若 $x\in B_{\frac{1}{2}\rho_2}(y)$，有 $F$ 的点 $z$ 位于 $B_\gamma(x)$ 内.由 \eqref{eq:4.9}，对任意 $\theta$，
\begin{equation}\label{eq:4.12}
  \alpha\gamma^{s}
  <\mathcal{H}^{s}\bigl(E\cap S_\gamma(z,\theta,\phi)\bigr)
  \leq\mathcal{H}^{s}\bigl(E\cap B_\gamma(z)\bigr)
  \leq\mathcal{H}^{s}\bigl(E\cap B_{2\gamma}(x)\bigr).
\end{equation}
若 $\gamma<\frac{1}{4}\rho_2$，则 $B_{\rho_2}(y)$ 含有 $\bigl(\rho_2/\gamma\bigr)^2/16$ 个圆心在 $B_{\frac{1}{2}\rho_2}(y)$ 中、半径为 $2\gamma$ 的不交圆盘.于是把这些圆盘上的 \eqref{eq:4.12} 求和，并由 \eqref{eq:4.10} 得
\[
  \bigl(\rho_2/\gamma\bigr)^{2}\alpha\gamma^{s}/16
  <\mathcal{H}^{s}\bigl(E\cap B_{\rho_2}(y)\bigr)
  <2^{s+1}\rho_2^{s},
\]
故 $\gamma>c\rho_2$，其中 $c$ 只依赖于 $\alpha$ 与 $s$.于是当 $\gamma\leq c\rho_2$ 时，存在含于 $B_{\rho_2}(y)$ 内且不含 $F$ 的点的半径为 $\gamma$ 的圆盘.从而可找到圆盘 $B_{\rho_3}(w)\subset B_{\rho_2}(y)$，其内部不含 $F$ 的点，但其边界含有 $F_0$ 的一点 $v$，且
\begin{equation}\label{eq:4.13}
  \rho_2\geq\rho_3\geq c\rho_2.
\end{equation}

\begin{figure}[htbp]
  \centering
  \includegraphics[width=0.36\textwidth]{fig4-3.png}
  \caption{\cref{thm:4.7} 证明中的圆盘与扇形}
  \label{fig:4.3}
\end{figure}

设 $\theta$ 为 $B_{\rho_3}(w)$ 在 $v$ 处的内法向方向，并令 $\rho_4$ 为 $B_{\rho_3}(w)$ 中过 $v$ 且与 $\theta$ 成角 $\phi$ 的弦长的一半（见图~\ref{fig:4.3}）.由初等几何，
\begin{equation}\label{eq:4.14}
  \rho_4=\rho_3\cos\phi.
\end{equation}
由于扇形 $S_{\rho_4}(v,\theta,\phi)$ 位于 $B_{\rho_3}(w)$ 内，它除 $v$ 外不含 $F$ 的点.于是由 \eqref{eq:4.11}、\eqref{eq:4.13} 与 \eqref{eq:4.14}（$c_1$ 只依赖于 $\phi,\alpha$ 与 $s$）得
\begin{equation}\label{eq:4.15}
  \mathcal{H}^{s}\bigl(E\cap S_{\rho_4}(v,\theta,\phi)\bigr)
  =\mathcal{H}^{s}\bigl((E\setminus F)\cap S_{\rho_4}(v,\theta,\phi)\bigr)
  \leq\mathcal{H}^{s}\bigl((E\setminus F)\cap B_{\rho_2}(y)\bigr)
  <\varepsilon\rho_2^{s}
  <\varepsilon c_1\rho_4^{s}.
\end{equation}
于是对任意 $\varepsilon>0$ 可找到 $v\in F_0$、满足 $0<\rho_4<\rho$ 的 $\rho_4$ 与 $\theta$ 使 \eqref{eq:4.15} 成立.这与 $F_0$ 的定义 \eqref{eq:4.9} 矛盾，故 $\mathcal{H}^s(F_0)=0$.注意到 $\alpha,\rho>0$ 任意，证明完成.
\end{proof}

\begin{corollary}\label{cor:4.8}
设 $E$ 是 $\mathbb{R}^2$ 中满足 $1<s<2$ 的 $s$-集.则在 $E$ 的几乎每一点处，下角密度 $D^s\bigl(E,x,\theta,\frac{1}{2}\pi\bigr)$ 对某个 $\theta$ 为零.
\end{corollary}

\begin{proof}
取递增趋于 $\frac{1}{2}\pi$ 的序列 $\{\phi_i\}$.由\cref{thm:4.7}，对几乎所有的 $x\in E$ 可找到方向序列 $\{\theta_i\}$ 使 $D^s(E,x,\theta_i,\phi_i)=0$ 对每个 $i$ 成立.抽取收敛子列，可设 $\theta_i\to\theta$.于是对任意 $\phi<\frac{1}{2}\pi$ 有 $D^s(E,x,\theta,\phi)=0$.因此，只要 $x$ 还不是\cref{lem:4.5} 的例外点，便有 $D^s\bigl(E,x,\theta,\frac{1}{2}\pi\bigr)=0$.
\end{proof}

下一步我们考察 $\mathbb{R}^2$ 中满足 $1<s<2$ 的正则 $s$-集的下角密度，作为证明这类集合不存在的一个步骤.\textcite{marstrand1954a} 的下述精巧论证大致上说：若这样的集合在过其一点的一条直线的一侧稀疏，则它在另一侧也必定稀疏.

\begin{lemma}\label{lem:4.9}
设 $E$ 是 $\mathbb{R}^2$ 中的 $s$-集，$1<s<2$.设 $x$ 是 $E$ 的上凸密度等于 $1$ 的正则点，并设对某个 $\theta$ 有 $D^s\bigl(E,x,-\theta,\frac{1}{2}\pi\bigr)=0$.则 $E$ 在 $x$ 处有垂直于 $\theta$ 的弱切线.
\end{lemma}

\begin{proof}
由于 $D^s(E,x)=1$ 且 $\bar{D}^s_{\mathrm{c}}(E,x)=1$，给定任意 $\eta>0$，可找到任意小的 $\rho$ 值使
\begin{equation}\label{eq:4.16}
  \mathcal{H}^{s}\bigl(E\cap B_r(x)\bigr)>2^{s}r^{s}(1-\eta)\qquad(r\leq\rho),
\end{equation}
\begin{equation}\label{eq:4.17}
  \mathcal{H}^{s}(E\cap U)<(1+\eta)|U|^{s}\qquad(x\in U\ \text{且}\ 0<|U|\leq 2\rho),
\end{equation}
并且
\begin{equation}\label{eq:4.18}
  \mathcal{H}^{s}\Bigl(E\cap S_\rho\bigl(x,-\theta,\tfrac{1}{2}\pi\bigr)\Bigr)<2^{s}\eta\rho^{s}.
\end{equation}
取 $0<\phi<\frac{1}{2}\pi$，设 $L$ 为过 $x$ 且垂直于 $\theta$ 的直线，$M$ 与 $M'$ 为从 $x$ 出发与 $\theta$ 成角 $\phi$ 的半射线.对固定的正整数 $m$，归纳地构造以 $x$ 为圆心、基于 $L$ 的 $m+1$ 个半径为 $r_i$ 的半圆 $S_{r_i}$，其中 $\rho=r_0>r_1>\cdots>r_m$.每个半圆 $S_{r_i}$ 的端点 $y_i$ 与 $y_i'$ 在 $L$ 上，且与 $M$ 和 $M'$ 分别交于 $z_i$ 与 $z_i'$.设 $S_{r_i}$ 已经构造，则如下指定 $S_{r_{i+1}}$：取 $y_{i+1}$ 为 $[x,y_i]$ 上使其到 $y_i'$ 的距离等于它到 $y_i$ 与到 $z_i'$ 的距离之和的那一点.（一个直截了当的连续性论证表明这样的点存在.）以 $y_{i+1}$ 为圆心、过 $z_i'$ 作弧交 $L$ 于 $y_{i+1}'$.由对称性，以 $y_{i+1}'$ 为圆心、过 $z_i$ 的弧交 $L$ 于 $y_{i+1}$（见图~\ref{fig:4.4}）.记 $S_{r_i}$ 被这两条弧截出的（阴影）凸部分为 $U_i$.则
\begin{equation}\label{eq:4.19}
  |U_i|=2r_{i+1}.
\end{equation}

\begin{figure}[htbp]
  \centering
  \includegraphics[width=0.44\textwidth]{fig4-4.png}
  \caption{\cref{lem:4.9} 证明中的半圆与区域 $U_i$}
  \label{fig:4.4}
\end{figure}

我们估计 $E$ 介于由 $M$ 与 $M'$ 界定的两个相邻扇形之间的部分的测度.由
\[
  S_{r_i}(x,\theta,\phi)\setminus S_{r_{i+1}}(x,\theta,\phi)
  \subset U_i\setminus S_{r_{i+1}}\bigl(x,\theta,\tfrac{1}{2}\pi\bigr)
  \subset U_i\cup S_\rho\bigl(x,-\theta,\tfrac{1}{2}\pi\bigr)\setminus B_{r_{i+1}}(x),
\]
得
\begin{align*}
  \mathcal{H}^{s}&\bigl(E\cap S_{r_i}(x,\theta,\phi)\bigr)-\mathcal{H}^{s}\bigl(E\cap S_{r_{i+1}}(x,\theta,\phi)\bigr)\\
  &\leq\mathcal{H}^{s}(E\cap U_i)+\mathcal{H}^{s}\Bigl(E\cap S_\rho\bigl(x,-\theta,\tfrac{1}{2}\pi\bigr)\Bigr)
    -\mathcal{H}^{s}\bigl(E\cap B_{r_{i+1}}(x)\bigr)\\
  &<(1+\eta)|U_i|^{s}+2^{s}\eta\rho^{s}-2^{s}r_{i+1}^{s}(1-\eta)\\
  &\leq 2^{s+1}\eta r_{i+1}^{s}+2^{s}\eta\rho^{s}
  \leq 2^{s+2}\eta\rho^{s},
\end{align*}
这里用了 \eqref{eq:4.16}--\eqref{eq:4.19}.把这一估计对全部 $m$ 个扇形求和，由 \eqref{eq:4.17} 得
\begin{align*}
  \mathcal{H}^{s}\bigl(E\cap S_\rho(x,\theta,\phi)\bigr)
  &<2^{s+2}\eta\rho^{s}m+\mathcal{H}^{s}\bigl(E\cap S_{r_m}(x,\theta,\phi)\bigr)\\
  &<2^{s+2}\eta\rho^{s}m+2^{s}(1+\eta)r_m^{s}.
\end{align*}
由构造，$r_m/\rho$ 只依赖于 $\phi$ 与 $m$，且当 $m\to\infty$ 时该比值趋于零.于是给定 $\varepsilon>0$，可找到与 $\eta$ 无关的 $m$ 使
\[
  \mathcal{H}^{s}\bigl(E\cap S_\rho(x,\theta,\phi)\bigr)<\rho^{s}\bigl(2^{s+2}\eta m+\varepsilon(1+\eta)\bigr).
\]
这对任意小的 $\rho$ 值与任意 $\eta>0$ 成立，故当 $\phi<\frac{1}{2}\pi$ 时 $D^s(E,x,\theta,\phi)=0$.由于 $D^s\bigl(E,x,-\theta,\frac{1}{2}\pi\bigr)=0$，由 \eqref{eq:4.7} 可得 $E$ 在 $x$ 处有垂直于 $\theta$ 的方向的弱切线.
\end{proof}

\begin{corollary}\label{cor:4.10}
设 $E$ 是 $\mathbb{R}^2$ 中满足 $1<s<2$ 的 $s$-集.则 $E$ 是非正则的.
\end{corollary}

\begin{proof}
由\cref{thm:2.3}，在 $E$ 的几乎所有的点处上凸密度等于 $1$；由\cref{cor:4.8}，对某个 $\theta$ 有 $D^s\bigl(E,x,\theta,\frac{1}{2}\pi\bigr)=0$.于是由\cref{lem:4.9}，集合 $E$ 在其几乎所有的正则点处有弱切线.这与\cref{cor:4.6} 不相容，除非正则点之集测度为零；此时 $E$ 是非正则 $s$-集.
\end{proof}

在发表上述结果后不久，\textcite{marstrand1955} 把该方法推广，证明了更强的事实：$\mathbb{R}^2$ 中（$1<s<2$）$s$-集的密度几乎处处不存在.

为完整起见，我们把本章的主要结果汇集在一起.

\begin{theorem}\label{thm:4.11}
平面 $s$-集除非 $s$ 是整数，必是非正则的.
\end{theorem}

\begin{proof}
合并\cref{cor:4.3} 与\cref{cor:4.10}.
\end{proof}

\section{高维集合}

主要的高维结果与\cref{thm:4.11} 类似.

\begin{theorem}\label{thm:4.12}
$\mathbb{R}^n$ 中的 $s$-集除非 $s$ 是整数，必是非正则的.
\end{theorem}

$0<s<1$ 的情形已在\cref{thm:4.2} 中证明，\cref{cor:4.10} 的证明稍作调整即可涵盖 $n-1<s<n$ 的情形.然而对中间的 $s$ 值，需要实质性的新思想，该定理最终由 \textcite{marstrand1964} 证明.事实上，Marstrand 的工作远不止于此：它不仅适用于更一般的测度，还证明了当 $s$ 非整数时密度几乎处处不存在.其证明包含对正则集的弱切向性质深入得多的研究.

\begin{exercises}
\item 调整\cref{thm:4.2} 的证明，证明若 $E$ 是满足 $0<s<1$ 的 $s$-集，则对几乎所有的 $x$ 有 $D^s(E,x)\leq\bigl(1+2^{s/(s-1)}\bigr)^{s-1}$.
\item 设 $E$ 是 $s$-集（$0<s<1$）.由\cref{thm:4.4} 推出：对任意单位向量 $\theta$，只要密度用开半圆或开半球计算，则在几乎所有的 $x\in E$ 处\newterm{density-hemispherical}{下半球密度}{density, hemispherical} $D^s\bigl(E,x,\theta,\frac{1}{2}\pi\bigr)=0$.举一个例子说明对闭半圆这 不成立.
\item 证明\cref{thm:4.12} 的如下特殊情形：设 $E$ 是 $\mathbb{R}^3$ 中位于光滑（比如说无穷可微）曲面上的 $s$-集.则除非 $s$ 是整数，$E$ 是非正则的.
\end{exercises}
